\documentclass{beamer}

\usepackage[brazilian]{babel}
\usepackage[T1]{fontenc}
\usepackage{graphicx}

% Tema
\usetheme{Metropolis}
\usecolortheme{default}

% Informações da apresentação
\title{Introdução ao LaTeX para Apresentações}
\subtitle{Criando slides com Beamer}
\author{Tiago Luis de Oliveira Moura}
\date{\today}

\begin{document}

% ==================================================
% CAPA
% ==================================================

\begin{frame}
    \titlepage
\end{frame}

% ==================================================
% SUMÁRIO
% ==================================================

\begin{frame}{Roteiro}
    \tableofcontents
\end{frame}

% ==================================================
\section{O que é LaTeX?}
% ==================================================

\begin{frame}{O que é LaTeX?}

    \begin{itemize}
        \item LaTeX é um sistema de preparação de documentos.
        \item É muito utilizado em artigos, teses, relatórios e apresentações científicas.
        \item O usuário descreve a estrutura do documento utilizando comandos.
        \item O LaTeX cuida da formatação e da composição final.
    \end{itemize}

\end{frame}

% ==================================================
\section{Beamer}
% ==================================================

\begin{frame}{O que é Beamer?}

    O \textbf{Beamer} é uma classe do LaTeX desenvolvida para a criação de apresentações.

    \vspace{0.5cm}

    Para utilizá-lo, iniciamos o documento com:

    \vspace{0.3cm}

    \begin{center}
        \texttt{\textbackslash documentclass\{beamer\}}
    \end{center}

\end{frame}

% ==================================================
\section{Estrutura básica}
% ==================================================

\begin{frame}{Estrutura básica}

    Uma apresentação Beamer possui uma estrutura semelhante a qualquer documento LaTeX.

    \vspace{0.4cm}

    \begin{block}{Código}
        \texttt{\textbackslash documentclass\{beamer\}}\\[0.2cm]

        \texttt{\textbackslash begin\{document\}}\\[0.2cm]

        \hspace{0.5cm}
        \texttt{\textbackslash begin\{frame\}\{Título\}}\\

        \hspace{1cm}
        \texttt{Conteúdo do slide}\\

        \hspace{0.5cm}
        \texttt{\textbackslash end\{frame\}}\\[0.2cm]

        \texttt{\textbackslash end\{document\}}
    \end{block}

\end{frame}

% ==================================================

\begin{frame}{O ambiente frame}

    Cada slide da apresentação é criado utilizando um ambiente
    \texttt{frame}.

    \vspace{0.4cm}

    \begin{block}{Exemplo}
        \texttt{\textbackslash begin\{frame\}\{Meu título\}}\\[0.2cm]

        \hspace{0.5cm}
        \texttt{Conteúdo do slide}\\[0.2cm]

        \texttt{\textbackslash end\{frame\}}
    \end{block}

\end{frame}

% ==================================================
\section{Texto e listas}
% ==================================================

\begin{frame}{Listas}

    Podemos organizar informações utilizando listas.

    \vspace{0.3cm}

    \begin{itemize}
        \item Primeiro item
        \item Segundo item
        \item Terceiro item
        \item Quarto item
    \end{itemize}

\end{frame}

% ==================================================

\begin{frame}{Listas numeradas}

    Também podemos utilizar listas numeradas:

    \begin{enumerate}
        \item Instalar uma distribuição LaTeX
        \item Criar um arquivo \texttt{.tex}
        \item Escrever o documento
        \item Compilar
        \item Gerar o PDF
    \end{enumerate}

\end{frame}

% ==================================================
\section{Equações}
% ==================================================

\begin{frame}{Equações}

    Uma das principais vantagens do LaTeX é a escrita de equações matemáticas.

    \vspace{0.5cm}

    \begin{equation}
        c^2 = a^2 + b^2
        \label{eq:pitagoras}
    \end{equation}

    \vspace{0.4cm}

    A Equação \ref{eq:pitagoras} representa o teorema de Pitágoras.

\end{frame}

% ==================================================

\begin{frame}{Código de uma equação}

    Para criar uma equação numerada:

    \vspace{0.4cm}

    \begin{block}{Código}
        \texttt{\textbackslash begin\{equation\}}\\

        \hspace{0.5cm}
        \texttt{c\^{}2 = a\^{}2 + b\^{}2}\\

        \texttt{\textbackslash end\{equation\}}
    \end{block}

\end{frame}

% ==================================================
\section{Imagens}
% ==================================================

\begin{frame}{Inserindo imagens}

    Imagens podem ser inseridas utilizando o comando
    \texttt{\textbackslash includegraphics}.

    \vspace{0.5cm}

    \begin{block}{Exemplo}
        \texttt{\textbackslash includegraphics}\\
        \texttt{[width=0.5\textbackslash linewidth]}\\
        \texttt{\{imagem.pdf\}}
    \end{block}

\end{frame}

% ==================================================
\section{Tabelas}
% ==================================================

\begin{frame}{Tabelas}

    \begin{center}

        \begin{tabular}{ccc}
            \hline
            A & B & C \\
            \hline
            1 & 2 & 3 \\
            4 & 5 & 6 \\
            7 & 8 & 9 \\
            \hline
        \end{tabular}

    \end{center}

\end{frame}

% ==================================================

\begin{frame}{Estrutura de uma tabela}

    \begin{block}{Código}

        \texttt{\textbackslash begin\{tabular\}\{ccc\}}\\[0.2cm]

        \hspace{0.5cm}
        \texttt{A \& B \& C \textbackslash\textbackslash}\\

        \hspace{0.5cm}
        \texttt{1 \& 2 \& 3 \textbackslash\textbackslash}\\

        \hspace{0.5cm}
        \texttt{4 \& 5 \& 6 \textbackslash\textbackslash}\\[0.2cm]

        \texttt{\textbackslash end\{tabular\}}

    \end{block}

\end{frame}

% ==================================================
\section{Links}
% ==================================================

\begin{frame}{Links}

    O Beamer já possui suporte para hyperlinks.

    \vspace{0.5cm}

    Um exemplo:

    \vspace{0.3cm}

    \begin{center}
        \href{https://www.tablesgenerator.com}{tablesgenerator.com}
    \end{center}

    \vspace{0.5cm}

    O código correspondente é:

    \begin{block}{Código}
        \texttt{\textbackslash href\{URL\}\{texto\}}
    \end{block}

\end{frame}

% ==================================================
\section{Temas}
% ==================================================

\begin{frame}{Temas do Beamer}

    O visual da apresentação pode ser modificado utilizando temas.

    \vspace{0.3cm}

    Esta apresentação utiliza:

    \begin{center}
        \texttt{\textbackslash usetheme\{Madrid\}}
    \end{center}

    Outros exemplos:

    \begin{itemize}
        \item Berlin
        \item CambridgeUS
        \item Copenhagen
        \item Warsaw
        \item Singapore
        \item Boadilla
    \end{itemize}

\end{frame}

% ==================================================
\section{Conclusão}
% ==================================================

\begin{frame}{Conclusão}

    \begin{itemize}
        \item Beamer é a classe do LaTeX utilizada para apresentações.
        \item Cada slide é construído com o ambiente \texttt{frame}.
        \item É possível inserir equações, tabelas, imagens e links.
        \item O conteúdo e a formatação ficam separados.
        \item É particularmente útil para apresentações acadêmicas e científicas.
    \end{itemize}

\end{frame}

% ==================================================

\begin{frame}
    \centering

    {\Huge Obrigado!}

    \vspace{0.8cm}

    {\Large Dúvidas?}
\end{frame}

\end{document}